\chapter{Questions worth being ready for}

\section{Why not just compare the two point clouds?}

Because absence in a point cloud is ambiguous. It can mean the object left, or it can mean the
camera never looked there. On real footage, where nobody walks the identical route twice, the
second case accounts for most of the room. Tracking confirmed empty space separates the two.

\section{Why one reconstruction instead of aligning two?}

Aligning two models means solving for the transform that makes them agree, and that is least
reliable exactly when part of the scene has moved, which is the only case this project cares
about. Doing it jointly lets the unchanged scene do the aligning, which is the right way round.

\section{How do you know it works?}

A synthetic room with changes scripted in advance. All four are found, in the right places, at
volumes consistent with the surface area of the objects. The figures in Chapter~5 are generated
from that run by code reading the output file, not drawn by hand.

\section{What are the weaknesses?}

Four, stated plainly.

\begin{itemize}
  \item Monocular depth is the weakest link, and its errors are smooth and plausible, which is the
        worst kind of error because nothing looks obviously wrong.
  \item There is no true metric scale, so volumes and distances are consistent but not in real
        metres.
  \item Six false positives appear on the test scene, from depth error on flat walls.
  \item The labelling has only been exercised on noise textured cubes, so it is untested on real
        objects.
\end{itemize}

\section{What would you do next?}

Real metric scale, from a known size marker in shot or from the phone's own motion tracking.
Dense stereo where a GPU is available. Log-odds occupancy, as used in OctoMap, so confidence
accumulates across frames instead of a hard threshold on how many times a cell was hit. An
open-vocabulary detector so objects outside the configured list can still be named.

\section{What did you get wrong along the way?}

Three things, all caught by testing rather than by reading the code.

\begin{enumerate}
  \item \textbf{The depth limit was written in metres while the reconstruction used arbitrary
        units.} The scene reconstructed at roughly twenty eight units deep, the limit was eight, so
        every single measurement was discarded as too far away and nothing fused at all. Fixed by
        rescaling the reconstruction once to room sized units.
  \item \textbf{The first synthetic room had no walls.} Any surface with a single corner behind the
        camera was being discarded whole, which removed the floor and every wall and left a black
        void. Fixed by cutting the large surfaces into tiles so the test applies locally.
  \item \textbf{The naming stage reported guesses as fact,} announcing that a television had moved
        when the model was scoring at chance level. Fixed with a confidence floor.
\end{enumerate}

\begin{plainbox}
All three were invisible from reading the code and obvious from running it against a known answer.
That is the argument for building the synthetic test room before trusting anything.
\end{plainbox}
